\section*{Hoofdstuk \ref{ch:g10:the-mole} --- De mol en de molaire massa}

\begin{solution}{exo:g10:the-mole:1}
$M(\ce{O2}) = 2 \times 16.0 = \qty{32.0}{g/mol}$;
$M(\ce{N2}) = 2 \times 14.0 = \qty{28.0}{g/mol}$;
$M(\ce{CH4}) = 12.0 + 4 \times 1.0 = \qty{16.0}{g/mol}$;
$M(\ce{NH3}) = 14.0 + 3 \times 1.0 = \qty{17.0}{g/mol}$;
$M(\ce{CO2}) = 12.0 + 2 \times 16.0 = \qty{44.0}{g/mol}$.
\end{solution}

\begin{solution}{exo:g10:the-mole:2}
Water: $n = 36.0 / 18.0 = \qty{2.00}{mol}$. Methaan:
$n = 4.0 / 16.0 = \qty{0.25}{mol}$.
\end{solution}

\begin{solution}{exo:g10:the-mole:3}
$m(\ce{NaCl}) = 0.25 \times 58.5 = \qty{14.6}{g}$, ongeveer \qty{15}{g};
$m(\ce{CO2}) = 2.0 \times 44.0 = \qty{88}{g}$.
\end{solution}

\begin{solution}{exo:g10:the-mole:4}
$N = 0.50 \times \num{6.02e23} = \num{3.0e23}$ zuurstofmoleculen.
$n = \num{3.0e24} / \qty{6.02e23}{mol^{-1}} = \qty{5.0}{mol}$ ijzer.
\end{solution}

\begin{solution}{exo:g10:the-mole:5}
$V = n \times V_m = 0.50 \times 24.1 = \qty{12}{L}$.
$n = V / V_m = 12 / 24.1 = \qty{0.50}{mol}$.
\end{solution}

\begin{solution}{exo:g10:the-mole:6}
De druppel heeft een massa van \qty{0.050}{g}, dus
$n = 0.050 / 18.0 = \qty{2.8e-3}{mol}$ en
$N = \num{2.8e-3} \times \num{6.02e23} = \num{1.7e21}$ \omterm{def:g7:atoms-and-molecules:molecule}{moleculen}.
\end{solution}

\begin{solution}{exo:g10:the-mole:7}
Een aluminiumatoom (\qty{27.0}{g/mol}) is ongeveer half zo zwaar als een
ijzeratoom (\qty{55.8}{g/mol}), dus dezelfde massa aluminium bevat
ongeveer twee keer zoveel \omterm{def:g7:atoms-and-molecules:atom}{atomen}. Controle: aluminium
$10/27.0 = \qty{0.370}{mol}$, dat is \num{2.23e23} \omterm{def:g7:atoms-and-molecules:atom}{atomen}; ijzer
$10/55.8 = \qty{0.179}{mol}$, dat is \num{1.08e23} \omterm{def:g7:atoms-and-molecules:atom}{atomen}. Aluminium wint,
met een factor van ongeveer 2.1.
\end{solution}

\begin{solution}{exo:g10:the-mole:8}
Eén liter is $1/24.1 = \qty{0.0415}{mol}$ gas. Koolstofdioxide:
$0.0415 \times 44.0 = \qty{1.83}{g}$; stikstof:
$0.0415 \times 28.0 = \qty{1.16}{g}$. Koolstofdioxide is ongeveer
anderhalf keer zo dicht als \omterm{def:g4:air-a-mixture-of-gases:air}{lucht}, waarvan de \omterm{def:g10:the-mole:molar-mass}{molaire massa} dicht bij die
van stikstof ligt. Het gistende sap geeft koolstofdioxide af, dat zakt en
zich ophoopt bij de vloer van een gesloten kelder, waar het iedereen die
naar binnen gaat kan doen stikken.
\end{solution}

\begin{solution}{exo:g10:the-mole:9}
$M(\ce{C6H12O6}) = 6 \times 12.0 + 12 \times 1.0 + 6 \times 16.0 =
\qty{180.0}{g/mol}$. Pijl $\div N_A$: $n = \num{1.0e22} / \num{6.02e23} =
\qty{1.66e-2}{mol}$; pijl $\times M$: $m = \num{1.66e-2} \times 180.0 =
\qty{3.0}{g}$.
\end{solution}

\begin{solution}{exo:g10:the-mole:10}
$M = \qty{342.0}{g/mol}$, dus $n = 5.0 / 342.0 = \qty{1.46e-2}{mol}$;
$N = \num{1.46e-2} \times \num{6.02e23} = \num{8.8e21}$ \omterm{def:g7:atoms-and-molecules:molecule}{moleculen}. Elk
heeft 12 koolstofatomen: $12 \times \num{8.8e21} = \num{1.1e23}$
koolstofatomen.
\end{solution}

\begin{solution}{exo:g10:the-mole:11}
Water: $1000 / 18.0 = \qty{55.6}{mol}$. Ethanol:
$M = 2 \times 12.0 + 6 \times 1.0 + 16.0 = \qty{46.0}{g/mol}$, dus
$1000 / 46.0 = \qty{21.7}{mol}$. Water heeft de grootste hoeveelheid, met
een factor $55.6 / 21.7 \approx 2.6$, wat gewoon $46.0 / 18.0$ is.
\end{solution}

\begin{solution}{exo:g10:the-mole:12}
Goud: $\qty{3.75}{g}$, $n = 3.75 / 197.0 = \qty{1.90e-2}{mol}$, dat is
\num{1.15e22} \omterm{def:g7:atoms-and-molecules:atom}{atomen}. Koper: \qty{1.25}{g},
$n = 1.25 / 63.5 = \qty{1.97e-2}{mol}$, dat is \num{1.19e22} \omterm{def:g7:atoms-and-molecules:atom}{atomen}.
Verrassend genoeg bevat de ring iets meer koperatomen dan goudatomen: een
goudatoom is ongeveer drie keer zo zwaar als een koperatoom.
\end{solution}

\begin{solution}{exo:g10:the-mole:13}
\begin{enumerate}
  \item $V = 8.0 \times 6.0 \times 3.0 = \qty{144}{m^3} =
        \qty{1.44e5}{L}$; $n = \num{1.44e5} / 24.1 = \qty{5.98e3}{mol}$;
        $N = \num{5.98e3} \times \num{6.02e23} = \num{3.6e27}$ \omterm{def:g7:atoms-and-molecules:molecule}{moleculen}.
  \item $M = (78 \times 28.0 + 21 \times 32.0 + 1 \times 40.0) / 100 =
        2896 / 100 = \qty{29.0}{g/mol}$.
  \item $m = \num{5.98e3} \times 29.0 = \qty{1.73e5}{g}$, ongeveer
        \qty{173}{kg}: de \omterm{def:g4:air-a-mixture-of-gases:air}{lucht} in een lokaal weegt evenveel als twee of
        drie volwassenen.
\end{enumerate}
\end{solution}

\begin{solution}{exo:g10:the-mole:14}
\begin{enumerate}
  \item $n = 1000 / 28.1 = \qty{35.6}{mol}$, dus
        $N = 35.6 \times \num{6.02e23} = \num{2.14e25}$ \omterm{def:g7:atoms-and-molecules:atom}{atomen}.
  \item $m = \qty{1.000}{kg} / \num{2.14e25} = \qty{4.67e-26}{kg}$.
  \item De massa van de bol en de \omterm{def:g10:the-mole:molar-mass}{molaire massa} geven de \omterm{def:g10:the-mole:amount}{hoeveelheid
        stof}, $n = m / M$. Als je het aantal \omterm{def:g7:atoms-and-molecules:atom}{atomen} $N$ onafhankelijk telt
        (uit het volume van de bol en de afstand tussen de \omterm{def:g7:atoms-and-molecules:atom}{atomen} in het
        kristal), is de verhouding $N / n$ de \omterm{def:g10:the-mole:avogadro}{constante van Avogadro}.
\end{enumerate}
\end{solution}

\begin{solution}{exo:g10:the-mole:15}
$M = 0.758 \times 35.0 + 0.242 \times 37.0 = 26.53 + 8.95 =
\qty{35.5}{g/mol}$: de waarde uit de tabel.
\end{solution}

\begin{solution}{pb:g10:the-mole:1}
\textbf{1.} $M(\ce{H2O}) = 2 \times 1.0 + 16.0 = \qty{18.0}{g/mol}$.

\textbf{2.} $m = 250 \times \qty{1.0}{g} = \qty{250}{g}$.

\textbf{3.} $n = 250 / 18.0 = \qty{13.9}{mol}$.

\textbf{4.} Elk \omterm{def:g7:atoms-and-molecules:molecule}{molecuul} heeft twee waterstofatomen en één zuurstofatoom:
\qty{27.8}{mol} waterstofatomen, \qty{13.9}{mol} zuurstofatomen.

\textbf{5.} $N = 13.9 \times \num{6.02e23} = \num{8.36e24}$ \omterm{def:g7:atoms-and-molecules:molecule}{moleculen}.

\textbf{6.} $m = \qty{18.0}{g/mol} / \qty{6.02e23}{mol^{-1}} =
\qty{2.99e-23}{g}$.

\textbf{7.} $\num{8.36e24} \times \qty{2.99e-23}{g} = \qty{250}{g}$, de
massa van het water in het glas, zoals het hoort.

\textbf{8.} $\num{8.36e24} / \num{3.16e7} = \num{2.6e17}$ jaar: een aantal
jaren met zeventien nullen, ver buiten elke mogelijke telling.

\textbf{9.} $\qty{1.335e9}{km^3} \times \num{e9} = \qty{1.335e18}{m^3}$,
en met $\qty{1}{m^3} = \qty{1000}{L}$, \qty{1.335e21}{L}.

\textbf{10.} $\num{8.36e24} / \num{1.335e21} = \num{6.26e3}$ gemarkeerde
\omterm{def:g7:atoms-and-molecules:molecule}{moleculen} per liter.

\textbf{11.} $\num{6.26e3} \times 0.250 = \num{1.57e3}$ gemarkeerde
\omterm{def:g7:atoms-and-molecules:molecule}{moleculen} in het tweede glas.

\textbf{12.} Massa $\num{1.335e21} \times \qty{1.0}{kg} =
\qty{1.335e21}{kg} = \qty{1.335e24}{g}$; $n = \num{1.335e24} / 18.0 =
\qty{7.42e22}{mol}$.

\textbf{13.} $13.9 / \num{7.42e22} = \num{1.87e-22}$.

\textbf{14.} $\num{1.87e-22} \times \num{8.36e24} = \num{1.57e3}$:
hetzelfde antwoord als bij vraag 11, zoals het hoort.

\textbf{15.} $\num{1.335e21} / 0.250 = \num{5.34e21}$ glazen.

\textbf{16.} $\num{8.36e24} / \num{5.34e21} \approx 1570$. Een glas bevat
ongeveer 1600 keer meer \omterm{def:g7:atoms-and-molecules:molecule}{moleculen} dan de oceanen glazen water bevatten;
gelijkmatig verspreid laten de \omterm{def:g7:atoms-and-molecules:molecule}{moleculen} van één glas er ongeveer 1600
achter in elk glas oceaanwater.

\textbf{17.} $1 / \num{6.26e3} = \qty{1.6e-4}{L} = \qty{0.16}{mL}$,
ongeveer drie druppels.

\textbf{18.} Het volume van de oceanen is een schatting die hoogstens tot
op drie of vier cijfers bekend is, en het werd afgerond; zeewater is geen
zuiver water (het bevat zouten); het glas is niet precies \qty{250}{mL};
en de menging door alle oceanen zou nooit volkomen gelijkmatig zijn.
Alleen de orde van grootte is betrouwbaar.

\textbf{19.} Ongeveer 1600 \omterm{def:g7:atoms-and-molecules:molecule}{moleculen} van het eerste glas komen terug in
het tweede.
\end{solution}
